% ╔══════════════════════════════════════════════════════════════════════════╗
% ║  EJERCICIOS RESUELTOS — LEYES DE LA LÓGICA PROPOSICIONAL                 ║
% ║  Curso: Aritmética | Semana: 01                                           ║
% ╚══════════════════════════════════════════════════════════════════════════╝

\newpage
\section{Ejercicios resueltos --- Leyes de la L\'{o}gica Proposicional}

\begin{tcolorbox}[
  enhanced, colback=GrisClaro, colframe=AzulPizarra,
  arc=2pt, boxrule=0.7pt, left=8pt, right=8pt, top=4pt, bottom=4pt
]
\small\sffamily
Las resoluciones presentan: identificación del método, análisis, desarrollo
paso a paso y respuesta final destacada. Se resuelven los ejercicios de
mayor representatividad por nivel.
\end{tcolorbox}
\vspace{0.4cm}

% ══ EJERCICIO 1 ══════════════════════════════════════════════════════════════
\begin{tcolorbox}[
  enhanced, breakable,
  colback=GrisClaro, colframe=AzulPizarra,
  fonttitle=\bfseries\sffamily,
  title={Ejercicio 1 --- Nivel B\'{a}sico},
  arc=3pt, boxrule=0.8pt
]
\begin{problema}
  Simplificar la expresión: $\;\sim(\sim p \vee \sim q)$
\end{problema}

\smallskip
\noindent
\textbf{\sffamily\color{AzulMedio}Tema:} De Morgan y doble negación
\hfill
\textbf{\sffamily\color{AzulMedio}M\'{e}todo:} Ley de De Morgan + Ley 1

\medskip
\noindent\textit{Observamos que} tenemos la negación de una disyunción.
Debemos aplicar De Morgan para ingresar la negación y luego verificar
si hay doble negación que simplificar.

\medskip
\noindent\textbf{\sffamily Desarrollo:}
\begin{align*}
  \sim(\sim p \vee \sim q)
  &\equiv \sim(\sim p) \wedge \sim(\sim q)
  && \text{De Morgan: $\sim(A \vee B) \equiv \sim A \wedge \sim B$} \\[6pt]
  &\equiv p \wedge q
  && \text{Doble negaci\'{o}n: $\sim(\sim p) \equiv p$}
\end{align*}

\noindent\textit{La negación exterior transforma la disyunción en conjunción
y cancela ambas negaciones internas. Es la aplicación directa combinada
de De Morgan (Ley 6) y la involutiva (Ley 1).}

\medskip
\begin{tcolorbox}[
  colback=AzulClaro!60, colframe=AzulPizarra,
  arc=2pt, boxrule=0.8pt, left=8pt, right=8pt
]
\centering
$\displaystyle \boxed{p \wedge q}$
\quad \textbf{Alternativa: C}
\end{tcolorbox}

\begin{cajatip}
  Cuando veas $\sim(\sim A \vee \sim B)$, aplica De Morgan y doble
  negación en un solo paso mental: el resultado siempre es $A \wedge B$.
  Identifica este patrón para ganar tiempo en el examen.
\end{cajatip}
\end{tcolorbox}
\vspace{0.5cm}

% ══ EJERCICIO 2 ══════════════════════════════════════════════════════════════
\begin{tcolorbox}[
  enhanced, breakable,
  colback=GrisClaro, colframe=AzulPizarra,
  fonttitle=\bfseries\sffamily,
  title={Ejercicio 2 --- Nivel B\'{a}sico (Absorci\'{o}n)},
  arc=3pt, boxrule=0.8pt
]
\begin{problema}
  Simplificar: $\;p \wedge (\sim p \vee q)$
\end{problema}

\smallskip
\noindent
\textbf{\sffamily\color{AzulMedio}Tema:} Absorción con negación
\hfill
\textbf{\sffamily\color{AzulMedio}M\'{e}todo:} Ley de absorción (Ley 9)

\medskip
\noindent\textit{Observamos que} la expresión tiene la forma
$p \wedge (\sim p \vee q)$, que corresponde exactamente a la tercera
forma de la ley de absorción.

\medskip
\noindent\textbf{\sffamily Desarrollo:}
\begin{align*}
  p \wedge (\sim p \vee q)
  &\equiv p \wedge q
  && \text{Absorci\'{o}n: } p \wedge (\sim p \vee q) \equiv p \wedge q
\end{align*}

\noindent\textit{Esta es una de las formas de absorción que más aparece
en exámenes. Reconocerla de forma directa evita desarrollos innecesarios.}

\medskip
\begin{tcolorbox}[
  colback=AzulClaro!60, colframe=AzulPizarra,
  arc=2pt, boxrule=0.8pt, left=8pt, right=8pt
]
\centering
$\displaystyle \boxed{p \wedge q}$
\quad \textbf{Alternativa: C}
\end{tcolorbox}

\begin{cajatip}
  Las cuatro formas de absorción con negación ($p \wedge (\sim p \vee q)$
  y $p \vee (\sim p \wedge q)$) deben memorizarse como patrones listos,
  pues aparecen frecuentemente en PUCP, UNI y San Marcos.
\end{cajatip}
\end{tcolorbox}
\vspace{0.5cm}

% ══ EJERCICIO 3 ══════════════════════════════════════════════════════════════
\begin{tcolorbox}[
  enhanced, breakable,
  colback=GrisClaro, colframe=AzulPizarra,
  fonttitle=\bfseries\sffamily,
  title={Ejercicio 3 --- Nivel Intermedio (PUCP)},
  arc=3pt, boxrule=0.8pt
]
\begin{problema}
  Determinar el equivalente más simple de:
  \[
    (p \vee q) \rightarrow (\sim p \wedge q)
  \]
\end{problema}

\smallskip
\noindent
\textbf{\sffamily\color{AzulMedio}Tema:} Condicional, De Morgan, distributiva
\hfill
\textbf{\sffamily\color{AzulMedio}M\'{e}todo:} Leyes 7, 6, 5, 10, 11

\medskip
\noindent\textit{Observamos que} la expresión es una condicional. El primer
paso es siempre eliminarla usando la Ley 7.

\medskip
\noindent\textbf{\sffamily Desarrollo:}
\begin{align*}
  &(p \vee q) \rightarrow (\sim p \wedge q) \\[6pt]
  &\equiv \sim(p \vee q) \vee (\sim p \wedge q)
  && \text{Ley 7: } A \to B \equiv \sim A \vee B \\[6pt]
  &\equiv (\sim p \wedge \sim q) \vee (\sim p \wedge q)
  && \text{De Morgan (Ley 6)} \\[6pt]
  &\equiv \sim p \wedge (\sim q \vee q)
  && \text{Distributiva (Ley 5)} \\[6pt]
  &\equiv \sim p \wedge \mathbf{V}
  && \text{Complemento (Ley 10)} \\[6pt]
  &\equiv \sim p
  && \text{Identidad (Ley 11)}
\end{align*}

\noindent\textit{El patrón clave está en el paso 3: al distribuir $\sim p$,
se genera $\sim q \vee q$ que por complemento es $\mathbf{V}$, y por
identidad se elimina, dejando únicamente $\sim p$.}

\medskip
\begin{tcolorbox}[
  colback=AzulClaro!60, colframe=AzulPizarra,
  arc=2pt, boxrule=0.8pt, left=8pt, right=8pt
]
\centering
$\displaystyle \boxed{\sim p}$
\quad \textbf{Alternativa: C}
\end{tcolorbox}

\begin{cajatip}
  Siempre que aparezca una condicional de la forma
  $(A \vee B) \rightarrow (\sim A \wedge B)$, el resultado simplificado
  será $\sim A$. Este patrón es recurrente en PUCP y Trilce.
\end{cajatip}
\end{tcolorbox}
\vspace{0.5cm}

% ══ EJERCICIO 4 ══════════════════════════════════════════════════════════════
\begin{tcolorbox}[
  enhanced, breakable,
  colback=GrisClaro, colframe=AzulPizarra,
  fonttitle=\bfseries\sffamily,
  title={Ejercicio 4 --- Nivel Avanzado (UNMSM)},
  arc=3pt, boxrule=0.8pt
]
\begin{problema}
  Determinar el esquema molecular de la proposición y simplificar:
  \textit{``Si el triángulo tiene dos lados iguales, entonces el triángulo
  se llama isósceles; y el triángulo no se llama isósceles; luego el
  triángulo no tiene dos lados iguales.''}
\end{problema}

\smallskip
\noindent
\textbf{\sffamily\color{AzulMedio}Tema:} Condicional, De Morgan, absorción
\hfill
\textbf{\sffamily\color{AzulMedio}M\'{e}todo:} Simbolización + Leyes 7, 9, 6, 10, 11

\medskip
\noindent\textbf{\sffamily Simbolización:}
\begin{itemize}
  \item $p$: El triángulo tiene dos lados iguales.
  \item $q$: El triángulo se llama isósceles.
\end{itemize}

\noindent Esquema molecular:
$\;[(p \rightarrow q) \wedge \sim q] \rightarrow \sim p$

\medskip
\noindent\textbf{\sffamily Desarrollo:}
\begin{align*}
  &[(p \rightarrow q) \wedge \sim q] \rightarrow \sim p \\[6pt]
  &\equiv [(\sim p \vee q) \wedge \sim q] \rightarrow \sim p
  && \text{Ley 7 (condicional)} \\[6pt]
  &\equiv (\sim q \wedge \sim p) \rightarrow \sim p
  && \text{Absorci\'{o}n: } (\sim p \vee q) \wedge \sim q
     \equiv \sim q \wedge \sim p \\[6pt]
  &\equiv \sim(\sim q \wedge \sim p) \vee \sim p
  && \text{Ley 7 (condicional)} \\[6pt]
  &\equiv (q \vee p) \vee \sim p
  && \text{De Morgan} \\[6pt]
  &\equiv q \vee (p \vee \sim p)
  && \text{Asociativa} \\[6pt]
  &\equiv q \vee \mathbf{V}
  && \text{Complemento (Ley 10)} \\[6pt]
  &\equiv \mathbf{V}
  && \text{Identidad (Ley 11)}
\end{align*}

\noindent\textit{La proposición completa es una \textbf{tautología}: es
verdadera independientemente del valor de $p$ y $q$. Esto tiene sentido
porque la estructura es la formalización del modus tollens, que es
siempre válido.}

\medskip
\begin{tcolorbox}[
  colback=AzulClaro!60, colframe=AzulPizarra,
  arc=2pt, boxrule=0.8pt, left=8pt, right=8pt
]
\centering
$\displaystyle \boxed{\mathbf{V} \quad \text{(Tautolog\'{i}a)}}$
\quad \textbf{Alternativa: D}
\end{tcolorbox}

\begin{cajatip}
  El esquema $[(p \rightarrow q) \wedge \sim q] \rightarrow \sim p$ es
  el \textbf{Modus Tollens}, que siempre es una tautología. Reconocer
  estructuras clásicas de inferencia válida ahorra tiempo en San Marcos
  y UNI.
\end{cajatip}
\end{tcolorbox}
\vspace{0.5cm}

% ══ EJERCICIO 5 ══════════════════════════════════════════════════════════════
\begin{tcolorbox}[
  enhanced, breakable,
  colback=GrisClaro, colframe=AzulPizarra,
  fonttitle=\bfseries\sffamily,
  title={Ejercicio 5 --- Nivel Avanzado: Operaci\'{o}n definida},
  arc=3pt, boxrule=0.8pt
]
\begin{problema}
  Se define la operación: $p \ast q \equiv (p \wedge \sim q) \vee (p \wedge q)$.

  Simplificar: $\;\sim[(p \ast \sim q) \rightarrow (\sim p \ast q)]$
\end{problema}

\smallskip
\noindent
\textbf{\sffamily\color{AzulMedio}Tema:} Operación definida, distributiva, condicional
\hfill
\textbf{\sffamily\color{AzulMedio}M\'{e}todo:} Reducción de la operación + Leyes 5, 7, 10

\medskip
\noindent\textbf{\sffamily Paso 1 — Reducir la operación $\ast$:}
\begin{align*}
  p \ast q
  &\equiv (p \wedge \sim q) \vee (p \wedge q) \\[4pt]
  &\equiv p \wedge (\sim q \vee q)
  && \text{Distributiva} \\[4pt]
  &\equiv p \wedge \mathbf{V}
  && \text{Complemento} \\[4pt]
  &\equiv p
  && \text{Identidad}
\end{align*}

\noindent\textit{La operación $p \ast q$ es siempre igual a $p$,
independientemente de $q$.}

\medskip
\noindent\textbf{\sffamily Paso 2 — Calcular los operandos:}
\begin{align*}
  p \ast \sim q &\equiv p
  && \text{(por lo demostrado arriba)} \\[4pt]
  \sim p \ast q &\equiv \sim p
  && \text{(el resultado es siempre el primer operando)}
\end{align*}

\medskip
\noindent\textbf{\sffamily Paso 3 — Sustituir y simplificar:}
\begin{align*}
  \sim[(p \ast \sim q) \rightarrow (\sim p \ast q)]
  &\equiv \sim[p \rightarrow \sim p]
  && \text{Sustituci\'{o}n} \\[6pt]
  &\equiv \sim[\sim p \vee \sim p]
  && \text{Ley 7 (condicional)} \\[6pt]
  &\equiv \sim[\sim p]
  && \text{Idempotencia (Ley 2)} \\[6pt]
  &\equiv p
  && \text{Doble negaci\'{o}n (Ley 1)}
\end{align*}

\medskip
\begin{tcolorbox}[
  colback=AzulClaro!60, colframe=AzulPizarra,
  arc=2pt, boxrule=0.8pt, left=8pt, right=8pt
]
\centering
$\displaystyle \boxed{p}$
\quad \textbf{Alternativa: E}
\end{tcolorbox}

\begin{cajatip}
  Cuando encuentres una operación definida, el primer paso siempre es
  simplificarla de forma general usando las leyes del álgebra. Luego
  sustiuye en la expresión principal. Nunca trabajes con la definición
  expandida dentro de una expresión compleja; eso lleva a errores.
\end{cajatip}
\end{tcolorbox}
\vspace{0.5cm}

% ══ EJERCICIO 6 ══════════════════════════════════════════════════════════════
\begin{tcolorbox}[
  enhanced, breakable,
  colback=GrisClaro, colframe=AzulPizarra,
  fonttitle=\bfseries\sffamily,
  title={Ejercicio 6 --- Circuito L\'{o}gico a Esquema Molecular (PUCP)},
  arc=3pt, boxrule=0.8pt
]
\begin{problema}
  Determinar el esquema molecular del siguiente circuito y reducirlo:
  dos ramas en paralelo; la rama superior tiene $p$ y $q$ en serie;
  la rama inferior tiene $\sim p$ y $q$ en serie. Todo ello en serie
  con un bloque de dos ramas en paralelo: $\sim p$ (superior) y $p$
  con $q$ en serie (inferior).
\end{problema}

\smallskip
\noindent
\textbf{\sffamily\color{AzulMedio}Tema:} Circuitos lógicos y simplificación
\hfill
\textbf{\sffamily\color{AzulMedio}M\'{e}todo:} Lectura de circuito + Leyes 9 y 5

\medskip
\noindent\textbf{\sffamily Paso 1 — Leer el circuito:}

\noindent El circuito corresponde al esquema molecular:
\[
  \bigl[(p \vee q) \wedge q\bigr] \wedge \bigl[\square p \vee (p \wedge q)\bigr]
\]

\noindent\textit{Nota: basado en el ejercicio 4 de la hoja de trabajo con
dos bloques en serie.}

\medskip
\noindent\textbf{\sffamily Paso 2 — Simplificar bloque por bloque:}
\begin{align*}
  \bigl[(p \vee q) \wedge q\bigr]
  &\equiv q
  && \text{Absorci\'{o}n: } q \wedge (p \vee q) \equiv q
\end{align*}

\begin{align*}
  q \wedge \bigl[\sim p \vee (p \wedge q)\bigr]
  &\equiv q \wedge \bigl[(\sim p \vee p) \wedge (\sim p \vee q)\bigr]
  && \text{Distributiva} \\[4pt]
  &\equiv q \wedge \bigl[\mathbf{V} \wedge (\sim p \vee q)\bigr]
  && \text{Complemento} \\[4pt]
  &\equiv q \wedge (\sim p \vee q)
  && \text{Identidad} \\[4pt]
  &\equiv q
  && \text{Absorci\'{o}n: } q \wedge (\sim p \vee q) \equiv q
\end{align*}

\medskip
\begin{tcolorbox}[
  colback=AzulClaro!60, colframe=AzulPizarra,
  arc=2pt, boxrule=0.8pt, left=8pt, right=8pt
]
\centering
$\displaystyle \boxed{q}$
\quad \textbf{Rpta.: $q$}
\end{tcolorbox}

\begin{cajatip}
  Al simplificar circuitos mixtos, trabaja bloque a bloque de izquierda
  a derecha. Identifica primero los bloques en paralelo (disyunción),
  luego conecta los bloques en serie (conjunción). La absorción casi
  siempre aparece en el paso final.
\end{cajatip}
\end{tcolorbox}

% FIN — resueltos.tex
